%  LaTeX support: latex@mdpi.com 
%  For support, please attach all files needed for compiling as well as the log file, and specify your operating system, LaTeX version, and LaTeX editor.

%=================================================================
\documentclass[information,article,submit,pdftex,moreauthors]{Definitions/mdpi} 

%=================================================================

% MDPI internal commands
\firstpage{1} 
\makeatletter 
\setcounter{page}{\@firstpage} 
\makeatother
\pubvolume{1}
\issuenum{1}
\articlenumber{0}
\pubyear{2023}
\copyrightyear{2023}
%\externaleditor{Academic Editor: Firstname Lastname}
\datereceived{7 March 2023} 
%\daterevised{} 
\dateaccepted{} 
\datepublished{} 
\hreflink{https://doi.org/} % If needed use \linebreak
%\doinum{}

%=================================================================
% Add packages and commands here. The following packages are loaded in our class file: fontenc, inputenc, calc, indentfirst, fancyhdr, graphicx, epstopdf, lastpage, ifthen, lineno, float, amsmath, setspace, enumitem, mathpazo, booktabs, titlesec, etoolbox, tabto, xcolor, soul, multirow, microtype, tikz, totcount, changepage, attrib, upgreek, cleveref, amsthm, hyphenat, natbib, hyperref, footmisc, url, geometry, newfloat, caption

\graphicspath{{figures/}} % path to folder with figures
\usepackage{subcaption}
\usepackage{listings}
\usepackage{xcolor}
\usepackage{gensymb} % for \textdegree
\usepackage{tabularx}
\usepackage{makecell}
\usepackage{array}
\newcolumntype{?}{!{\vrule width 1pt}}

%=================================================================
% Full title of the paper (Capitalized)
%\Title{Land Features Recognition in Qena Bend of the Nile River, Egypt Using Unsupervised Classification of the Landsat Images by R}
\Title{Recognizing the Wadi Fluvial Structure and Stream Network in Qena Bend of the Nile River, Egypt on Landsat 8-9 OLI Images}

% MDPI internal command: Title for citation in the left column
\TitleCitation{Recognizing the Wadi Fluvial Structure and Stream Network in Qena Bend of the Nile River, Egypt on Landsat 8-9 OLI Images}

% Author Orchid ID: enter ID or remove command
\newcommand{\orcidauthorA}{0000-0002-5759-1089} % Polina Add \orcidA{} behind the author's name
\newcommand{\orcidauthorB}{0000-0002-6461-1551} % Olivier Add \orcidB{} behind the author's name

% Authors, for the paper (add full first names)
\Author{Polina Lemenkova $^{1}$*\orcidA{} and Olivier Debeir $^{1}$\orcidB{}}

%\longauthorlist{yes}

% MDPI internal command: Authors, for metadata in PDF
\AuthorNames{Polina Lemenkova and Olivier Debeir}

% MDPI internal command: Authors, for citation in the left column
\AuthorCitation{Lemenkova, P.; Debeir, O.}
% If this is a Chicago style journal: Lastname, Firstname, Firstname Lastname, and Firstname Lastname.

% Affiliations / Addresses (Add [1] after \address if there is only one affiliation.)
\address{%
$^{1}$ \quad Laboratory of Image Synthesis and Analysis (LISA), École polytechnique de Bruxelles (Brussels Faculty of Engineering), Université Libre de Bruxelles (ULB). Building L, Campus du Solbosch, ULB -- LISA CP165/57, Avenue Franklin D. Roosevelt 50, Brussels 1050, Belgium.
}

% Contact information of the corresponding author
\corres{Correspondence: polina.lemenkova@ulb.be; Tel.: +32 471 86 04 59 (P.L.)}

% Abstract
\abstract{With remote sensing data processing methods widely prevalent, the ability to handle changes in spatial data and evaluate distribution of diverse landforms across target areas in the datasets becomes important. One way to approach this problem is through satellite image processing. In this paper, we primarily focus on the methods of the unsupervised classification of the Landsat OLI/TIRS images covering the region of Qena governorate in Upper Egypt. We present a case of extracting information from the space-borne data for analysis of distribution of wadi streams over several years: 2013, 2015, 2016, 2019, 2022, and 2023. Six images were processed using computer vision methods of R. The results of k-means clustering of each scene were compared to visualise changes in landforms retrieved from the images covering Qena Bend of the Nile River. }

% Keywords
\keyword{Africa; k-means clustering; image processing; remote sensing; unsupervised classification; programming; visualization; modelling; computer vision; cartography} 

\PACS{91.10.Da; 91.10.Jf; 91.10.Sp; 91.10.Xa; 96.25.Vt; 91.10.Fc; 95.40.+s; 95.75.Qr; 95.75.Rs; 42.68.Wt}
\MSC{86A30; 86-08; 86A99; 86A04}
\JEL{Y91; Q20; Q24; Q23; Q3; Q01; R11; O44; O13; Q5; Q51; Q55; N57; C6; C61}

%%%%%%%%%%%%%%%%%%%%%%%%%%%%%%%%%%%%%%%%%%
\begin{document}

\section{Introduction}

Satellite images are frequently used as a source of geoinformation for mapping land cover types and recognising land features. Example of such applications include thematic maps made using classification of satellite images; these include, for instance maps of mineral resources \cite{ABDELKAREEM2018682,1294518,rs11121450}, environmental assessment \cite{615837,Lemenkova202227}, land cover/land use classification \cite{516791,Lemenkova202229}, geologic analysis \cite{6723607,KAMEL2022104420} and urban mapping \cite{MOHAMED2019103507}. The success of the remote sensing data in cartographic processing and mapping is based on the effectiveness of the satellite imagery as a source of information for recognition of diverse land features and processes. 

Among the advantages of the high-resolution space borne data as a source of information one can mention the improved quality of spatial visualization \cite{ABDALLA20128,Lemenkova202230}, and the enabled access to the remotely located places and areas otherwise inaccessible, such as deserts \cite{6350376}. For such places, using satellite images presents a valuable source of information. Apart from spatial access, satellite images provide the temporal analysis of the environmental changes or disaster monitoring using comparative analysis of several scenes. For instance, the comparison of several satellite images covering target region in different time periods enables to perform time series analysis which can be used for detecting climate and environmental changes, monitoring disaster events and assess the risks of hazards, and to perform the ecological surveillance of the target study area using multi-temporal data analysis \cite{9171725,SINGH2020423,Hashim}. 

Applications of the remote sensing data in Earth sciences include monitoring the desertification and land degradation \cite{Badreldin,Kawy}, analysis of deforestation \cite{516791}, evaluation of salinization \cite{8949893} and soil erosion \cite{Gad2020}, hydrological modelling and quantification of the flooded areas \cite{TAHA2017157,ELSADEK2019377} and many more. Spectral patterns of the multispectral Landsat images are especially widely used in geosciences as a source of spatial information for land cover mapping and recognition of land features \cite{8486245}. To this end, the geometric, radiometric, and spectral precision of the Landsat scenes can be enhanced through a combination of the high-resolution panchromatic and multispectral channels by fusion methods \cite{4420687}.

Natural resource monitoring benefits from the use of satellite images which provide a powerful source of information for detecting the environmental processes \cite{Kamel,Khafagy}. Further, identifying lithologic land surface structures and extracting extracted structural lineaments is possible using computer vision techniques \cite{Bishta,Kusky,Salman}. For instance, research on mineral exploration can be supported using the classified satellite images to discriminate lithological units on the colour composites of the images in various ratios \cite{Shalaby}. For geomorphological modelling, the Shuttle Radar Topography Mission (SRTM) Digital Elevation Model (DEM) can be overlaid with satellite images \cite{Liping}. Such data integration enables to identify the morphometric parameters of the river catchment and drainage area, and to examine spatial distribution of the agriculture fields with regards to the fluvial networks \cite{6265340}. 

Information derived from the classified satellite images enables to perform the assessment of geomorphological and geological risks for defining maps of hazards. This includes the problem of morphometric modelling using remote sensing data for identification of geologic hazards, as widely investigated in cartography, geoinformatics and Earth science communities. Hamdan and Khozyem \cite{Hamdan} applied the ASTER data and GIS techniques to delineate drainage networks in Wadi El-Mathula watershed in Central Eastern Desert. Various methods of geomorphometric modelling exist and can be applied to detect the spatial patterns of fluvial systems, or for hydrogeological modelling aimed at the analysis of groundwater potential \cite{Issawi,Abdalla2020,Gaber}. The geometry and morphology of the relief, the characteristics of topographic parameters such as slope, aspect and hillshade can be extracted as morphometric characteristics for analysis of watershed using 3D mapping \cite{Lemenkova202301}. 

Extracting the information from the remotely sensed data can be used for complex hydrological analysis, for instance, predicting wadi flash floods using morphometric parameters that represent the topographic and drainage characteristics of the basins using the geographic information system (GIS) \cite{w9070553}. Further, such information can be used for the analysis of wadi networks to define basin boundaries and drainage networks, to evaluate stream channel density, and to model slope and flow directions. Moreover, it is useful for retrospective modelling of the hydrological processes in the past to better analyse geologic setting in present using robust remote sensing data. Such approach enables deeper insights in the structure of the land surface visible from the space borne data \cite{abdelkareem2013integration,AbdelkareemFarouk}. 

Further, the essential geographic information collected from the satellite imagery is useful for operative flash flood monitoring through integration of the data on physiographic features. Other applications of the remote sensing data include modelling lithological, geological, and structural features \cite{BESHR2021999}, mapping mineral resources \cite{Kamal} and detecting wadi channels based on the combination of the topographic maps and satellite images \cite{Moubark}. 

In this paper, we propose an R-based case of the satellite image processing to address the problem of computer vision application for for extracting geoinformation from the remotely sensed space borne data. The main motivation of this work, is to demonstrate that land surface types can be extracted from a short time series of multispectral Landsat 8-9 OLI/TIRS images by k-means clustering by libraries of R language. This work reports an empirical comparison of the classification of the six Landsat scenes for years 2013, 2015, 2016, 2019, 2022 and 2023 to detect changes in the identified land surface types for target study region on Qena Bend, Upper Egypt. Furthermore, the results report a comparison of the detected classes by Landsat bands in false colour composite (5-4-3 channels corresponding to the Near-Infrared, Red and Green Bands in the Landsat OLI'TIRS images). The results on the performed classification support the argument that satellite images can be integrated out to achieve a recognition of the land surface types in a complex drainage network of Wadi Qena, Upper Egypt, and this paper demonstrates the practical approach to achieve satellite image processing by scripts of R programming language.

%\pagebreak
\section{Prior Work}

Qena Bend is one of the most notable morphological features of the Nile River in Upper Egypt which has a special characteristic curvature in its geometry (Figure \ref{fig01}). 

\begin{figure}[H]
\centering
	\includegraphics[width=14.0 cm]{Fig_01.jpg}
	\caption{Topographic map of Egypt. Target study area of Qena Bend region is shown in red rotated square. Software: GMT v. 6.1.1. Data source: GEBCO/SRTM. Cartography source: authors.
	\label{fig01}}
\end{figure}

The structural geomorphology of the Qena Bend largely reflects the general geologic setting of the Nile River and major processes of its evolution \cite{Said1981,said2017geology}. The Valley of the Nile River has been cut during late Miocene time of the Neogene period and formed later on as a results of the regional and local tectonics. The discovery of paleorivers in the Sahara by radar imaging systems revieled the orientation of the paleochannels and their flow directions \cite{ABDELKAREEM201235}. Major geomorphic structure of the Nile valley follows the faults oriented parallel to the Red Sea along the central current of the Nile and the Gulf of Suez in the north \cite{Ali2022,Youssef2003StructuralSO} with the relics of the pre-late Miocene streams remaining in fluviatile sand and gravel sediments of the Nile delta in Lower Egypt \cite{Said1981a}. 

The Qena Formation includes the gravel deposits with limestone cobble and heavily weathered soil typical for old terrace deposits \cite{paulissen1987egyptian}. The remaining sediments, outcrops and paleospecies in fossils are identified from the Mesozoic Era (Cretaceous period) \cite{ABDELHAMID2014138,SOLIMAN1986111} and Paleogene (Paleocene and Eocene epoches) \cite{MANDUR2022104594,sehim1993cretaceous}. They are reflected in the Cenozoic sediments accumulated in the lacustrine-alluvial fans around the Qena Bend due to the Nile River evolution in Neogene. During Neogene, it is argued \cite{PHILOBBOS2015194} that the ancient Qena Lake broke up the Nile River course which might have affected its present special form and specific geometric curvature (Figure \ref{fig02}). 

\begin{figure}[H]
\centering
	\includegraphics[width=13.0 cm]{Fig_02}
	\caption{Enlarged fragment of the 3D perspective view of Qena Bend, Nile River, Upper Egypt. Software: GMT v. 6.1.1. Data source: GEBCO/SRTM. Cartography source: authors.
	\label{fig02}}
\end{figure}

As a result of the geologic evolution, the geomorphology of the Qena Bend region present a plateau largely dissected by numerous valleys of wadi and geologic lineaments which dissect the cliffs of the Nile River and basement plateau composed of the limestone and chalks \cite{BANDEL1987427}. This result in a specific geomorphic pattern of the Qena Bend surroundings which is notable for a complex drainage network of wadi aquifer systems \cite{Alkholy,AbdelMoneim,HUSSIEN201773}. Such drainage network presents a typical feature of the Eastern Desert which consist of numerous fluvial basins that drain the rainwater towards the Nile River or the Red Sea \cite{Abdel}.

Apart from the geology and hydrology of the Nile River, the region of Qena Bend is influenced by the presence of the Eastern Desert which has a special climate. Its special feature includes the occasional extreme flash floods -- one of the most severe natural disasters. Flash floods occur annually, mostly during spring and autumn, and are caused by the torrential rainfalls followed by the disastrous surface runoff thereafter \cite{Moawad}. Such climate processes, aggravated by an interrelated complex of fluvial networks in Wadi Qena catchment area, as briefly explained above, presents favourable conditions for disastrous flash floods. During such events, a complex of the ephemeral rivers dry during most of the year is being filled with water from heavy showers and rains \cite{ElFakharany}. 

The cumulative effects from wadi network and arid and semi-arid climate of the Eastern Desert region results in occasional flash flood events and irregular sand dust storms with different durations and intensities. Flash floods occurs in diverse regions across in Egypt including Upper Egypt with Eastern Desert \cite{ElHaddad,Elsadek}, southern Red Sea Coast \cite{Abdalla2014}, Sinai Peninsula \cite{ELOSTA2021101522} and Gulf of Suez \cite{AbdelLattif}. The social consequences of flash floods include damaged infrastructure, transport ways and buildings, and occasional victims \cite{Moawad}. El-Magd et al. report the catastrophic events of flash floods in Qena governorate recorded in three consecutive events in 2014--2016 \cite{ElMagd}. 

The occurrence of flood hazards in Upper Egypt point at the significance of the application of geoinformation for visualization and mapping the regions prone to flood risk. The remote sensing data and advanced methods of spatial data processing present robust approaches to evaluation of relief and extracting information on relief and topographic characteristics from the satellite images using methods of computer vision. Such approaches can be used, for instance for vulnerability mapping or identification of areas with high geomorphological hazards and expose to floods \cite{nhess-9-751-2009}.

%%%%%%%%%%%%%%%%%%%%%%%%%%%%%%%%%%%%%%%%%%
\section{Materials and Methods}

\subsection{Data}

Multispectral satellite images such as Landsat TM/ETM+, and OLI/TIRS) provide a valuable source of remote sensing data widely used in environmental studies of Egypt \cite{Elfadaly,Sultanetal1986,Sultanetal1987}. In this work, we use the two major types of data Figure \ref{fig02}: the topographic grid for mapping the study area and a series of six satellite images presented by the Landsat 8-9 OLI/TIRS scenes. The high-resolution topographic data were used for inspection of the relief in the region, while the remote sensing data were used for comparative analysis of changes in the wadi channels in the Qena Bend area and surroundings by image processing and k-means clustering algorithms using R \cite{RCoreTeam}. 

\begin{figure}[H]
\centering
	\includegraphics[width=14.0 cm]{Fig_03}
	\caption{Data sources: summary of the materials used in this study. Software: R version 4.2.2, 'Mermaid' library. Flowchart source: authors.
	\label{fig03}}
\end{figure}

The detailed information on the satellite images is summarised in Table \ref{tab1}. 

\begin{table}[H]
\small
\caption{Satellite images used for computing the VIs: Landsat-8 OLI/TIRS collected from the USGS. \textsuperscript{1}.\label{tab1}}
	\begin{adjustwidth}{-\extralength}{0cm}
%		\newcolumntype{C}{>{\centering\arraybackslash}X}
		\begin{tabularx}{\fulllength}{ p{1.6cm} | p{1.8cm} | p{7.0cm} | p{4.0cm} | p{2.0cm} }
			\toprule
			\textbf{Date} & \textbf{Spacecraft} & \textbf{Landsat Product ID} & \textbf{Scene ID} & \textbf{Cloudiness} \\
			\midrule
			2013/11/16 & Landsat 8 & \makecell{LC08\_L2SP\_175042\_20131116\_20200912\_02\_T1} & LC81750422013320LGN01 & 0.89 \\
			2015/11/22 & Landsat 8 & \makecell{LC08\_L2SP\_175042\_20151122\_20200908\_02\_T1} & LC81750422015326LGN01 & 0.05 \\
			2016/11/08 & Landsat 8 & \makecell{LC08\_L2SP\_175042\_20161108\_20200905\_02\_T1} & LC81750422016313LGN01 & 0.31 \\
			2019/11/17 & Landsat 8 & \makecell{LC08\_L2SP\_175042\_20191117\_20200825\_02\_T1} & LC81750422019321LGN00 & 1.16 \\
			2022/11/09 & Landsat 8 & \makecell{LC08\_L2SP\_175042\_20221109\_20221121\_02\_T1} & LC81750422022313LGN00 & 0.02 \\
			2023/01/20 & Landsat 9 & \makecell{LC09\_L2SP\_175042\_20230120\_20230122\_02\_T1} & LC91750422023020LGN01 & 0.32 \\	
			\bottomrule
		\end{tabularx}
	\end{adjustwidth}
	\noindent{\footnotesize{\textsuperscript{1} The Sensor ID is common for all the scenes: Landsat 8-9 OLI/TIRS (Operational Land Imager and Thermal Infrared Sensor), Collection 2 Level-2. Path/row parameters common for all the images: 175/42. Coverage: Qena Bend, the Nile River, Upper Egypt. Image source: the USGS}}
\end{table}

\begin{figure}[H]
\makebox[1.00\linewidth][r]{%
	\begin{subfigure}[b]{.40\textwidth}
		\centering
			\includegraphics[width=0.87\textwidth]{Fig_04a.jpg}
		\caption{2013}
	\end{subfigure}%
	\begin{subfigure}[b]{.40\textwidth}
		\centering
		\includegraphics[width=0.87\textwidth]{Fig_04b.jpg}
		\caption{2015}
	\end{subfigure}%
	\begin{subfigure}[b]{.40\textwidth}
		\centering
		\includegraphics[width=0.87\textwidth]{Fig_04c.jpg}
		\caption{2018}
	\end{subfigure}%
}\\
\vfill \vspace{1mm}
\makebox[1.00\linewidth][r]{%
	\begin{subfigure}[b]{.40\textwidth}
		\centering
			\includegraphics[width=0.87\textwidth]{Fig_04d.jpg}
		\caption{2020}
	\end{subfigure}%
	\begin{subfigure}[b]{.40\textwidth}
		\centering
		\includegraphics[width=0.87\textwidth]{Fig_04e.jpg}
		\caption{2021}
	\end{subfigure}%
	\begin{subfigure}[b]{.40\textwidth}
		\centering
		\includegraphics[width=0.87\textwidth]{Fig_04f.jpg}
		\caption{2022}
	\end{subfigure}%
}\caption{Qena Bend of the Nile River, Egypt visible on Landsat 8-9 OLI/TIRS images in natural RGB colour for six years: (\textbf{a}) 2013/12/20, (\textbf{b}) 2015/12/26, (\textbf{c}) 2018/02/21, (\textbf{d}) 2020/01/22, (\textbf{e}) 2021/12/18, (\textbf{f}) 2022/12/29.}\label{fig:04}
\end{figure}

%----------- subsection ----------->
\subsection{Methods}

The 2D and 3D maps showing the topography and geomorphic structure of the Wadi Qena region were prepared using the Generic Mapping Tools \cite{Wessel} by existing methods of cartographic scripting \cite{Lemenkova202228}.

All the maps derived from the satellite images were georeferenced automatically by R library 'raster' in the WGS84 datum, UTM Coordinate system Zone 36, 

\begin{figure}[H]
\centering
	\includegraphics[width=12.0 cm]{Fig_05}
	\caption{Methodological workflow of the processes and research steps applied in this study. Software: R version 4.2.2, 'DiagrammeR' library. Flowchart source: authors.
	\label{fig05}}
\end{figure}

%%%%%%%%%%%%%%%%%%%%%%%%%%%%%%%%%%%%%%%%%%
\pagebreak
\section{Results}

The system of Wadi Fluvial structure and drainage stream network in the surroundings of the Qena Bend of the Nile River is largely controlled by the underlying geological structure and regional lithology. The origin of these basins takes place in mountainous highland of the basement and is oriented West-East, either to the Nile or to the Red Sea, according to catchment. In both cases, the potential damage from the disastrous events may affect urban areas of the small cities and infrastructure of the industrial places. Figure \ref{fig06} shows the wadi streams with their typical dendrite structure, oriented towards Qena Bend segment of the Nile.

\begin{figure}[H]
\makebox[1.00\linewidth][r]{%
	\begin{subfigure}[b]{.40\textwidth}
		\centering
			\includegraphics[width=0.87\textwidth]{Fig_06a}
		\caption{2013: Bands 432}
	\end{subfigure}%
	\begin{subfigure}[b]{.40\textwidth}
		\centering
		\includegraphics[width=0.87\textwidth]{Fig_06b}
		\caption{2013: Bands 543}
	\end{subfigure}%
	\begin{subfigure}[b]{.40\textwidth}
		\centering
		\includegraphics[width=1.00\textwidth]{Fig_06c}
		\caption{2013: Clustering}
	\end{subfigure}%
}\\
\vfill \vspace{1mm}
\makebox[1.00\linewidth][r]{%
	\begin{subfigure}[b]{.40\textwidth}
		\centering
			\includegraphics[width=0.87\textwidth]{Fig_06d}
		\caption{2015: Bands 571}
	\end{subfigure}%
	\begin{subfigure}[b]{.40\textwidth}
		\centering
		\includegraphics[width=0.87\textwidth]{Fig_06e}
		\caption{2015: Bands 632}
	\end{subfigure}%
	\begin{subfigure}[b]{.40\textwidth}
		\centering
		\includegraphics[width=0.92\textwidth]{Fig_06f}
		\caption{2015: Clustering}
	\end{subfigure}%
}\caption{Clustering of the Landsat 8-9 OLI/TIRS images: Qena Bend of the Nile River visible on in natural RGB colours for six years: (\textbf{a}) 2013 RGB in 432 Bands, (\textbf{b}) 2013 RGB in 543 Bands, (\textbf{c}) 2013 unsupervised classification, (\textbf{d}) 2015 RGB in 571 Bands, (\textbf{e}) 2015 RGB in 632 Bands, (\textbf{f}) 2015 RGB: Clustering.}\label{fig06}
\end{figure}

\begin{table}[H]
\small
\caption{Results of the k-means classification of the Landsat-8 OLI images for Qena Bend area\textsuperscript{1}.\label{tab2}}
	\begin{adjustwidth}{-\extralength}{0cm}
		\begin{tabularx}{\fulllength}{ C | C | C | C | C | C | C }
			\toprule
			\textbf{Class ID} & \textbf{2013} & \textbf{2015} & \textbf{2016} & \textbf{2019} & \textbf{2022} & \textbf{2023} \\
			\midrule
			1 & 6237758.6 & 25058427 & 11830859 & 11050520 & 33773365.2 & 22414794 \\
			2 & 26530972.6 & 3829403 & 6847520 & 7687698 & 21467288.5 & 4080038 \\
			3 & 46356528.0 & 23160871 & 13056437 & 31925587 & 466472.5 & 13918989 \\
			4 & 10314256.8 & 7458695 & 12594525 & 3359878 & 3547592.0 & 1863680 \\
			5 & 7861914.8 & 9909609 & 3012708 & 11101016 & 7714730.1 & 6083803 \\
			6 & 625035.3 & 4568955 & 2195275 & 14290386 & 9346294.4 & 14743313 \\
			7 & 22694195.3 & 3019950 & 12277811 & 14898378 & 6978872.8 & 6370890 \\
			8 & 5804917.1 & 1942973 & 10144633 & 5386949 & 16906933.2 & 6053497 \\
			9 & 7320502.4 & 13450745 & 3220846 & 5501193 & 8429786.2 & 11952344 \\
			10 & 6121676.5 & 10228437 & 4195904 & 10594134 & 11161798.4 & 10517903 \\	
			\bottomrule
		\end{tabularx}
	\end{adjustwidth}
	\noindent{\footnotesize{\textsuperscript{1} Sum of squares by clusters within computed cluster centroids.}}
\end{table}

\begin{figure}[H]
\makebox[1.00\linewidth][r]{%
	\begin{subfigure}[b]{.40\textwidth}
		\centering
			\includegraphics[width=0.87\textwidth]{Fig_07a}
		\caption{2016: Bands 564}
	\end{subfigure}%
	\begin{subfigure}[b]{.40\textwidth}
		\centering
		\includegraphics[width=0.87\textwidth]{Fig_07b}
		\caption{2016: Bands 754}
	\end{subfigure}%
	\begin{subfigure}[b]{.40\textwidth}
		\centering
		\includegraphics[width=0.90\textwidth]{Fig_07c}
		\caption{2016: Clustering}
	\end{subfigure}%
}\\
\vfill \vspace{1mm}
\makebox[1.00\linewidth][r]{%
	\begin{subfigure}[b]{.40\textwidth}
		\centering
			\includegraphics[width=0.87\textwidth]{Fig_07d}
		\caption{2019: Bands 764}
	\end{subfigure}%
	\begin{subfigure}[b]{.40\textwidth}
		\centering
		\includegraphics[width=0.87\textwidth]{Fig_07e}
		\caption{2019: Bands 753}
	\end{subfigure}%
	\begin{subfigure}[b]{.40\textwidth}
		\centering
		\includegraphics[width=0.92\textwidth]{Fig_07f}
		\caption{2019: Clustering}
	\end{subfigure}%
}\caption{Clustering of the Landsat 8-9 OLI/TIRS images Qena Bend of the Nile River, Egypt visible on in natural RGB colour for six years: (\textbf{a}) 2016 RGB in 564 Bands, (\textbf{b}) 2016 RGB in 754 Bands, (\textbf{c}) 2016 unsupervised classification, (\textbf{d}) 2019 RGB in 764 Bands, (\textbf{e}) 2019 RGB in 753 Bands, (\textbf{f}) 2019 RGB: Clustering.}\label{fig:07}
\end{figure}

\begin{table}[H]
\small
\caption{Pairwise comparison of the k-means classification for Bands 5-4-3 of the Landsat-8 OLI/TIRS images for 2013 and 2015 in Qena Bend region. \textsuperscript{1}.\label{tab2}}
	\begin{adjustwidth}{-\extralength}{0cm}
		\begin{tabularx}{\fulllength}{||C?C|C|C?C|C|C||}
			\toprule
			\multirow{2}{*}{$Class ID$} & \multicolumn{3}{c}{\textbf{2013}} & \multicolumn{3}{c}{\textbf{2015}} \\ 
			\textbf{} & \emph{Band 5} & \emph{Band 4} & \emph{Band 3} & \emph{Band 5} & \emph{Band 4} & \emph{Band 3} \\
			\midrule
			1 & 21019.08 & 18348.92 & 15722.58 & 19819.50 & 17226.143 & 14765.14 \\
			2 & 21956.38 & 19148.12 & 16294.50 & 12087.00 & 10917.333 & 10115.00 \\
			3 & 26892.38 & 23105.00 & 19120.92 & 18736.17 & 9956.333 & 10142.17 \\
			4 & 14308.50 & 13163.75 & 11730.75 & 24835.46 & 21496.385 & 17990.38 \\
			5 & 19246.50 & 16942.69 & 14626.81 & 23315.56 & 20178.688 & 17043.31 \\
			6 & 17656.18 & 15405.09 & 13460.00 & 26302.78 & 22807.222 & 19156.67 \\
			7 & 23465.45 & 20280.18 & 17179.82 & 15028.50 & 13507.500 & 12109.33 \\
			8 & 16549.67 & 10211.17 & 10132.83 & 29412.00 & 25219.500 & 20731.50 \\
			9 & 22484.00 & 19481.25 & 16802.50 & 17737.46 & 15833.385 & 13783.23 \\
			10 & 24895.53 & 21740.20 & 18326.47 & 17737.46 & 18747.389 & 15990.33 \\	
			\bottomrule
		\end{tabularx}
	\end{adjustwidth}
	\noindent{\footnotesize{\textsuperscript{1} Clusters are computed for each band in the composite triplets: Band 5 (NIR), Band 4 (Red) and Band 3 (Green)}}
\end{table}

\begin{table}[H]
\small
\caption{Pairwise comparison of the k-means classification for Bands 5-4-3 of the Landsat-8 OLI/TIRS images for 2016 and 2019 in Qena Bend region. \textsuperscript{1}.\label{tab2}}
	\begin{adjustwidth}{-\extralength}{0cm}
		\begin{tabularx}{\fulllength}{||C?C|C|C?C|C|C||}
			\toprule
			\multirow{2}{*}{$Class ID$} & \multicolumn{3}{c}{\textbf{2016}} & \multicolumn{3}{c}{\textbf{2019}} \\
			\textbf{} & \emph{Band 5} & \emph{Band 4} & \emph{Band 3} & \emph{Band 5} & \emph{Band 4} & \emph{Band 3} \\
			\midrule
			1 & 25067.53 & 21805.82 & 18283.65 & 20721.42 & 18438.42 & 15725.33 \\
			2 & 19837.27 & 17506.45 & 15060.45 & 24198.84 & 21005.26 & 17763.84 \\
			3 & 27148.44 & 23483.89 & 19596.67 & 25536.93 & 22109.73 & 18481.20 \\
			4 & 16049.43 & 13237.00 & 11880.43 & 18257.00 & 9376.75 & 9707.50 \\
			5 & 19541.25 & 9323.50 & 9742.75 & 14476.67 & 12599.83 & 11411.17 \\
			6 & 12780.50 & 11767.00 & 10670.50 & 16992.00 & 15236.88 & 13304.75 \\
			7 & 21370.47 & 18891.65 & 16208.94 & 18829.53 & 16928.47 & 14654.00 \\
			8 & 23122.89 & 20069.28 & 16927.33 & 31113.00 & 26355.00 & 21077.00 \\
			9 & 18240.88 & 16536.62 & 14205.25 & 27920.00 & 23667.00 & 19115.00 \\
			10 & 17269.14 & 15186.43 & 13284.43 & 22736.25 & 19858.25 & 16938.06 \\	
			\bottomrule
		\end{tabularx}
	\end{adjustwidth}
	\noindent{\footnotesize{\textsuperscript{1} Clusters are computed for each band in the composite triplets: Band 5 (NIR), Band 4 (Red) and Band 3 (Green)}}
\end{table}

\pagebreak

\begin{figure}[H]
\makebox[1.00\linewidth][r]{%
	\begin{subfigure}[b]{.40\textwidth}
		\centering
			\includegraphics[width=0.87\textwidth]{Fig_08a}
		\caption{2022: Bands 631}
	\end{subfigure}%
	\begin{subfigure}[b]{.40\textwidth}
		\centering
		\includegraphics[width=0.87\textwidth]{Fig_08b}
		\caption{2022: Bands 543}
	\end{subfigure}%
	\begin{subfigure}[b]{.40\textwidth}
		\centering
		\includegraphics[width=0.90\textwidth]{Fig_08c}
		\caption{2022: Clustering}
	\end{subfigure}%
}\\
\vfill \vspace{1mm}
\makebox[1.00\linewidth][r]{%
	\begin{subfigure}[b]{.40\textwidth}
		\centering
			\includegraphics[width=0.87\textwidth]{Fig_08d}
		\caption{2023: Bands 562}
	\end{subfigure}%
	\begin{subfigure}[b]{.40\textwidth}
		\centering
		\includegraphics[width=0.87\textwidth]{Fig_08e}
		\caption{2023: Bands 753}
	\end{subfigure}%
	\begin{subfigure}[b]{.40\textwidth}
		\centering
		\includegraphics[width=0.92\textwidth]{Fig_08f}
		\caption{2023: Clustering}
	\end{subfigure}%
}\caption{Clustering of the Landsat 8-9 OLI/TIRS images Qena Bend of the Nile River, Egypt visible on in natural RGB colour for six years: (\textbf{a}) 2022 RGB in 631 Bands, (\textbf{b}) 2022 RGB in 543 Bands, (\textbf{c}) 2022 unsupervised classification, (\textbf{d}) 2023 RGB in 562 Bands, (\textbf{e}) 2023 RGB in 753 Bands, (\textbf{f}) 2023 RGB: Clustering.}\label{fig:08}
\end{figure}

\begin{table}[H]
\small
\caption{Pairwise comparison of the k-means classification for Bands 5-4-3 of the Landsat-8 OLI/TIRS images for 2022 and 2023 in Qena Bend region. \textsuperscript{1}.\label{tab2}}
	\begin{adjustwidth}{-\extralength}{0cm}
		\begin{tabularx}{\fulllength}{||C?C|C|C?C|C|C||}
			\toprule
			\multirow{2}{*}{$Class ID$} & \multicolumn{3}{c}{\textbf{2022}} & \multicolumn{3}{c}{\textbf{2023}} \\ 
			\textbf{} & \emph{Band 5} & \emph{Band 4} & \emph{Band 3} & \emph{Band 5} & \emph{Band 4} & \emph{Band 3} \\
			\midrule
			\midrule
			1 & 24428.50 & 21422.20 & 18051.50 & 21513.50 & 10245.750 & 10483.000 \\
			2 & 19238.75 & 17313.00 & 14911.25 & 25528.40 & 22234.800 & 18665.200 \\
			3 & 27402.67 & 23631.83 & 19472.50 & 17918.11 & 9096.444 & 9575.333 \\
			4 & 22303.50 & 19676.83 & 16793.33 & 24549.75 & 21362.375 & 17997.875 \\
			5 & 17517.33 & 11594.67 & 10989.33 & 15965.40 & 14835.200 & 13574.400 \\
			6 & 12943.67 & 11050.00 & 10303.00 & 22865.67 & 19456.933 & 16537.733 \\
			7 & 17225.22 & 15122.22 & 13253.00 & 23488.23 & 20695.000 & 17726.000 \\
			8 & 25640.82 & 22323.82 & 18730.73 & 18460.75 & 16531.833 & 14379.417 \\
			9 & 20699.44 & 18365.81 & 15840.12 & 20155.57 & 18063.786 & 15478.143 \\
			10 & 23487.22 & 20428.72 & 17181.22 & 27425.50 & 23626.900 & 19359.200 \\	
			\bottomrule
		\end{tabularx}
	\end{adjustwidth}
	\noindent{\footnotesize{\textsuperscript{1} Clusters are computed for each band in the composite triplets: Band 5 (NIR), Band 4 (Red) and Band 3 (Green)}}
\end{table}

%%%%%%%%%%%%%%%%%%%%%%%%%%%%%%%%%%%%%%%%%%
\pagebreak
\section{Discussion}

The Qena Wadi region is affected by annual flash floods and related geomorphological processes such as erosion as a consequence of the soil mass movement and deposition of soil particles during the heavy rainfalls. In turn, the deterioration of soil and changes in land cover types negatively affect the agriculture and may damage infrastructure, e.g., broken roads. Such consequences of the natural hazards well illustrate the dependence of social sustainability and urban infrastructure on the environmental setting of the region. Accurate information regarding the possibility of geomorphic risks or landslide deposition may be derived from evaluation of the landscapes's morphology. For instance, the evaluation of the channel incision modelled by geomorphic adjustment scenarios if an example of the effects from geomorphic processes on infrastructure \cite{Hermas}.

In our study, we exploited the time series of the satellite images for classification, topographic maps (GEBCO/SRTM grids) for 2D and 3D mapping for illustration of the geomorphology and topography of the Wadi Qena region, and statistical data derived from the image analysis for evaluating the change in land cover types aimed at assessing the geomorphological hazard risks in Upper Egypt. The integration of such data as well as analysis of the geological, geomorphological and climate setting based on prior works was used to evaluate the variations in land cover types illustrating the risks of hazards in the Qena Bend area. Similar examples of data integration for hydrogeological modelling are presents in earlier studies focused on evaluating the aquifers potentiality in Eastern Desert \cite{Moneim,ELHAMEED2017218} 2D and 3D modelling for land surface slope gradient assessment \cite{Othman}, recognising the structural lineaments and fluvial channels and catchment analysis \cite{rs12101559}. Such studies the the approaches of data fusion which profits from the advanced computer vision algorithms for extracting the information from remote sensing data \cite{AZZAZY2022104805}.

Computer vision algorithms of image recognition, discrimination and data processing are fundamental problem in remote sensing and Earth sciences and has a wide range of applications \cite{ABDELKAREEM2015245,Lemenkova202231,ELQADI2020104696,Husseinetal1995,DHANYA2022211}. Visualization, analysis and classification of satellite images by computer vision algorithms aims at preventing the risk of natural hazards through mapping of the potentially endangered areas. Such regions as wadi valleys in desert areas of southern Egypt are at high risks of flash floods. With this regard, this paper presents an approach to represent, process, analysis and compare satellite images, both as classified maps and as extracted information on land cover types and their changes as statistical data derived from the images and summarised in tables. We show how an R-based approach based on k-means clustering and spectral reflectance characteristics of the remotely sensed data can be used to learn how landscapes in southern region of Egypt gradually change over time. 

The discrimination of the images was performed to recognise these changes using a short-time series analysis of several images using principles of remote sensing data processing \cite{curran1985principles}. The analysis of similarity in texture and structure of the objects represented on the Landsat 8-9 scenes was performed based on the characteristics of the multispectral channels from the Landsat satellite images using fusion techniques for divers spectral bands \cite{4234117}. The combination of the false colour composites (NIR, Red and Green) was applied to better recognise the attributes of vegetation, water and land as major classes by automatic machine-based algorithms of R. Consequently, the subclasses such as wadi valleys, Nile floodplain and valley, hills, were defined according to the surface structure in the study area. As can be recognised on the satellite images, general orientation of the patterns of Wadi Qena coincides almost parallel to the axis direction of the Red Sea uplift, which corresponds with the geologic origin of the wadi network.

Flash floods and soil deposition in Wadi Qena region are climate-related natural phenomena with repetitive occurrence typical for southern Egypt due to specific geomorphology \cite{ElShamy}. Therefore, predicting such hazards is a subject of complex climatic computational modelling. Nevertheless, using robust information derived from the remote sensing data processed by advanced tools of computer vision and pattern recognition enables to support measures on risk mitigation and propose planning strategies in regions prone to natural hazards. Moreover, geomorphic data can be used for evaluating the groundwater suppliance \cite{Abdelkhalek2021,Abdelkareem} and its their availability for irrigation in agriculture \cite{Kamal2021}. Therefore, the satellite-derived land cover types maps are a useful tools for determining hazard areas and defining the zones susceptible to risks of floods and applied geomorphic modelling according to the distribution of fluvial network of wadi in southern segment of the Eastern Desert. Moreover, data analysis and visualization performed using scripting methods provide a new source of information that can be used for taking the mitigation policies and measures of sustainable development in Qena Band region, as well as urban planning in hazard prone areas of Egypt \cite{ABDELAAL2020103671}. 

In the future works, the assessment of the geomorphological hazards by analysis of susceptibility of various regions supported by the classified satellite images can be further used for the analysis of the flash flood, soil erosion and and mass deposition for evaluation of hazard risks. Since changes in landscapes are subject to destructive geomorphological hazards such as flash floods, dust storms and heavy showers, using satellite images presents a source of objective and robust information regarding these hazards and disastrous processes. Because frequently occurring processes of flash floods cause deposition of sediments and debris and affect present shape of the land surface, the preventive mapping and evaluation of the geomorphological risks is important both for infrastructure and agricultural activities in Upper Egypt. 

%%%%%%%%%%%%%%%%%%%%%%%%%%%%%%%%%%%%%%%%%%
\section{Conclusions}
In this paper, we have demonstrated a method to extract information  from the satellite images Landsat 8-9 OLI/TIRS to track land cover changes in wadi drainage surface of Qena Bend district of the Nile River. We have presented an application of R scripting libraries for processing geoinformation to analyse visual objects on the remote sensing data by computer vision and methods of pattern recognition on the images. Local features in the target study area were identified using k-means clustering method and visualized on maps based on the classified six images. This algorithm exploits the values of the pixels' DNs to discriminate the classes on the land surface based on the similarity of pixels' values and assigning them to 10 classes. 

We have approached the problem of unsupervised classification of the satellite images by R for extracting information on the land cover types as objects in Qena Bend region of the Upper Egypt. Hence, the objects representing spatial features and land cover classes were separated automatically in context of geomorphological structures over the studied territory using machine-based approach by R programming. The application of R libraries 'raster', 'terra' and 'RStoolbox' enabled to retrieve information form the remote sensing data and analyse changes in the complex of the ephemeral riverbeds in Qena district by methods of computer vision and applied programming.

In the framework of this study, the scripting implementation of the R libraries has been evaluated under Landsat OLI/TIRS data with auxiliary cartographic data processing by GMT. The comparison of the land cover changes has been summarised in the tables based on the information extracted from the R-processed images. Our analyses shows that R approach and algorithms of computer vision applied for geographic data processing is effective for multispectral space borne data. The R-based method proposes an advanced approach to geospatial image processing compared to state of the art GIS methods. We demonstrated the effectiveness of R in automatic interpretation of satellite images based on algorithm trained to classify pixels and assign them to classes according to spectral reflectance of the land cover classes in various bands of the satellite images. 

Due to high level of automation and application of computer vision algorithms, the approach of computer vision was used to present a series of several maps illustrating the land cover types in Qena area the periods of 2013, 2015, 2016, 2019, 2022 and 2023.  The comparison of changes in land cover types helps to better understand and interpret the cumulative effects from the environmental changes and the effects from the occasional floods in southern segment of the Eastern Desert on the geomorphic patterns in land cover types of Qena Bend. 

%%%%%%%%%%%%%%%%%%%%%%%%%%%%%%%%%%%%%%%%%%
\authorcontributions{Supervision, conceptualization, methodology, software, resources, funding acquisition, and project administration, O.D.; writing---original draft preparation, methodology, software, data curation, visualization, formal analysis, validation, writing---review and editing, and investigation, P.L. All authors have read and agreed to the published version of the manuscript.}

\funding{The publication was funded by the Editorial Office of \emph{Information}, MDPI, by providing 100\% discount for the APC of this manuscript. This project was supported by the Federal Public Planning Service Science Policy or Belgian Federal Science Policy Office - BELSPO (B2/202/P2/SEISMOSTORM).}

\institutionalreview{Not applicable.}

\informedconsent{Not applicable.}

\dataavailability{Not applicable} 

\acknowledgments{The authors thank the reviewers for reading and review of this manuscript.}

\conflictsofinterest{The authors declare no conflict of interest.} 

\abbreviations{Abbreviations}{
The following abbreviations are used in this manuscript:\\

\noindent 
\begin{tabular}{@{}ll}
MDPI & Multidisciplinary Digital Publishing Institute\\
DOAJ & Directory of open access journals\\
LD & Linear dichroism
\end{tabular}
}

\begin{adjustwidth}{-\extralength}{0cm}
%\printendnotes[custom] % Un-comment to print a list of endnotes

\bibliography{Qena.bib}
\reftitle{References}
\nocite{*}

%=====================================
% References, variant A: external bibliography
%=====================================


%%%%%%%%%%%%%%%%%%%%%%%%%%%%%%%%%%%%%%%%%%
\end{adjustwidth}
\end{document}
